\section{Scalar geometry of the Lubetzky--Zhao boundary}
\label{app:lz-scalar-geometry}

This appendix gives the proofs of the three supporting lemmas used in
\Cref{sec:lz-boundary}, together with the detailed verifications of (M1), (M3),
and (M4) in \Cref{thm:scalar-lz-boundary}.  They are collected here to
separate the scalar convex-analysis details from the construction of the
boundary functions $\pc$ and $\sm$ in the main text.

We shall use the following standard two-point representation of a convex
minorant.

\begin{lemma}[Two-point representation of convex minorants]
\label{lem:two-point-convex-minorant}
Let $I\subset\mathbb R$ be a compact interval, let $f:I\to\mathbb R$ be
continuous, and let $\widehat f$ be the convex minorant of $f$.  Then, for
every $x\in I$,
\[
  \widehat f(x)
  =
  \min\left\{
    \Big(\lambda f(x_1)+(1-\lambda)f(x_2)\Big):
      x_1,x_2\in I,\  \lambda\in[0,1],\
      \lambda x_1+(1-\lambda)x_2=x
  \right\}.
\]
\end{lemma}
\begin{proof}
The general representation theorem expresses the convex minorant as an
infimum over convex combinations of values of $f$.  In one dimension, two
points suffice, and compactness of $I$ and continuity of $f$ ensure that the
infimum is attained.  See, for example,
\cite[Corollary~17.1.5]{RockafellarConvexAnalysis}.
\end{proof}

\begin{proof}[Proof of \Cref{lem:convexity-defect}]
Differentiating the definition of $h_{p,d}$ and using
\[
  J_p'(z)=\log\frac{z(1-p)}{(1-z)p},
  \qquad
  J_p''(z)=\frac1{z(1-z)},
  \qquad
  J_p'''(z)=\frac{2z-1}{z^2(1-z)^2}
\]
gives
\[
  h_{p,d}'(z)=(2-d)J_p''(z)+zJ_p'''(z)
  =\frac{d(z-r_*)}{z(1-z)^2}.
\]
The denominator is positive.  Hence, $h_{p,d}'$ is negative before $r_*$ and
positive after $r_*$, which proves that $h_{p,d}$ is strictly decreasing on
$(0,r_*)$, strictly increasing on $(r_*,1)$, and has its unique minimum at
$r_*$.

We next evaluate this minimum.  Since $r_*=(d-1)/d$,
\[
  r_*J_p''(r_*)=d,
  \qquad
  J_p'(r_*)=\log\frac{(d-1)(1-p)}{p},
\]
and therefore
\[
  h_{p,d}(r_*)
  =d-(d-1)\log\frac{(d-1)(1-p)}{p}.
\]
As a function of $p$, this value is strictly increasing because its derivative
is
\[
  (d-1)\left(\frac1{1-p}+\frac1p\right)>0.
\]
Solving $h_{p,d}(r_*)=0$ gives
\[
  p=\frac{d-1}{(d-1)+\exp(d/(d-1))}=p_*.
\]
It follows that $h_{p,d}(r_*)\ge0$ exactly when $p\ge p_*$.  Since $r_*$ is
the minimum point, $h_{p,d}\ge0$ throughout $(0,1)$ in this case.

Now suppose that $p<p_*$.  Then, $h_{p,d}(r_*)<0$.  The alternative expression
\[
  h_{p,d}(z)=\frac1{1-z}-(d-1)J_p'(z)
\]
shows that $h_{p,d}(z)\to+\infty$ as $z\downarrow0$ and as $z\uparrow1$.
At the left endpoint $J_p'(z)\to-\infty$, while at the right endpoint
$1/(1-z)$ dominates the logarithmic divergence of $J_p'(z)$.  The intermediate
value theorem and the strict monotonicity on either side of $r_*$ therefore
give unique zeros $u_-(p)<r_*<u_+(p)$.  The same monotonicity gives the stated
sign of $h_{p,d}$ on the three intervening intervals.

It remains to translate this sign information into curvature of
$\varphi_{p,d}$.  For $z=x^{1/d}$, the chain rule gives
\[
  \varphi_{p,d}'(x)=\frac{J_p'(z)}{d z^{d-1}},
  \qquad
  \varphi_{p,d}''(x)
  =\frac{zJ_p''(z)-(d-1)J_p'(z)}{d^2z^{2d-1}}
  =\frac{h_{p,d}(z)}{d^2z^{2d-1}}.
\]
The denominator is positive, so $\varphi_{p,d}''(x)$ has the same sign as
$h_{p,d}(z)$.  When $p\ge p_*$, this proves convexity on $(0,1)$, and
continuity extends it to $[0,1]$.  When $p<p_*$, the change of variables
$x=z^d$ converts the three sign intervals for $h_{p,d}$ into precisely the
convexity and concavity intervals stated in the lemma.
\end{proof}

\begin{proof}[Proof of \Cref{lem:contact-points}]
Write $\varphi=\varphi_{p,d}$ and $u_\pm=u_\pm(p)$, put
\[
  \alpha=u_-^d,
  \qquad
  \beta=u_+^d,
\]
and let $\widehat\varphi$ be the convex minorant of $\varphi$.  By
Lemma~\ref{lem:convexity-defect}, $\varphi$ is strictly convex on
$(0,\alpha)$ and $(\beta,1)$ and strictly concave on $(\alpha,\beta)$.
The proof has three parts.  First we use the convex minorant to produce the
two contact points and identify the shape of the minorant.  Then, we prove that
no other pair of contact points can satisfy the same common-tangent equations.
Finally we record the analytic dependence on $p$ and the monotonicity of the
two contacts.

We begin with one elementary tangency observation: if an affine
function $L$ lies below a differentiable function $\varphi$ in a neighborhood
of an interior point $t$, and if $L(t)=\varphi(t)$, then the slope of $L$ is
$\varphi'(t)$.

We first extract the two contacts from the convex minorant.  The convex
minorant cannot touch $\varphi$ at any point $c\in(\alpha,\beta)$.  To see
this, suppose $\widehat\varphi(c)=\varphi(c)$ and take a supporting line to the
convex function $\widehat\varphi$ at $c$.  This line lies below
$\widehat\varphi$, hence below
$\varphi$, and has equality with $\varphi$ at $c$.  By the tangency
observation, this line has slope $\varphi'(c)$.  But $\varphi$ is strictly
concave near $c$, and the tangent line to a strictly concave function lies
strictly above the graph on both sides of the tangency point.  This is a
contradiction.  Thus,
\[
  (\alpha,\beta)\subset\{x : \widehat\varphi(x)<\varphi(x)\}.
\]

Let $(a,b)$ be the connected component of
$\{x : \widehat\varphi(x)<\varphi(x)\}$ containing $(\alpha,\beta)$.  Fix
$x\in(a,b)$.  By Lemma~\ref{lem:two-point-convex-minorant}, there are
$s,t\in[0,1]$ and $\lambda\in[0,1]$ such that
\[
  x=\lambda s+(1-\lambda)t,
  \qquad
  \widehat\varphi(x)=\lambda\varphi(s)+(1-\lambda)\varphi(t).
\]
After interchanging $s$ and $t$ and replacing $\lambda$ by $1-\lambda$ if
necessary, we must have
$s<x<t$ and $0<\lambda<1$; otherwise the value in this representation would
be $\varphi(x)$.  Convexity of $\widehat\varphi$ and the inequality
$\widehat\varphi\le\varphi$ give
\[
  \widehat\varphi(x)
  \le \lambda\widehat\varphi(s)+(1-\lambda)\widehat\varphi(t)
  \le \lambda\varphi(s)+(1-\lambda)\varphi(t)
  =\widehat\varphi(x).
\]
Thus, both inequalities are equalities.  Since both weights are positive,
$\widehat\varphi(s)=\varphi(s)$ and
$\widehat\varphi(t)=\varphi(t)$.  For a convex function, equality at a
nontrivial convex combination forces equality with the chord on the whole
interval; hence $\widehat\varphi$ is affine on $[s,t]$.
Since $(a,b)$ contains no contact points, we must have $s\le a$ and $t\ge b$.
Hence
$\widehat\varphi=L$ on $(a,b)$ for some affine function $L$.

The endpoints of this component are contact points.  At an interior
endpoint this follows from the continuity of
$\varphi-\widehat\varphi$.  At the endpoints of the domain,
Lemma~\ref{lem:two-point-convex-minorant} gives
\[
  \widehat\varphi(0)=\varphi(0),
  \qquad
  \widehat\varphi(1)=\varphi(1),
\]
because the only convex combination of points of $[0,1]$ with barycenter
$0$ (respectively, $1$) is supported at $0$ (respectively, $1$).
Consequently, the affine function $L$ extends continuously to $[a,b]$, and
\[
  L(a)=\varphi(a),\qquad L(b)=\varphi(b),
\]
where the value of $L$ at a domain endpoint is its corresponding one-sided
limit from $(a,b)$.  We next rule out these endpoint cases.

First suppose $a=0$.  Since $L$ is affine, it has some finite slope $m$, and
the contact condition at the endpoint gives $L(0)=\varphi(0)$.  The derivative
of $\varphi$ satisfies $\varphi'(x)\to-\infty$ as $x\downarrow0$.  By the mean
value theorem, for all sufficiently small $x>0$,
\[
  \frac{\varphi(x)-\varphi(0)}{x}<m.
\]
Using $L(x)=L(0)+mx=\varphi(0)+mx$, this inequality becomes
\[
  \varphi(x)<\varphi(0)+mx=L(x),
\]
contradicting $L=\widehat\varphi\le\varphi$ on $(a,b)$.

Since $\varphi'(x)\to+\infty$ as $x\uparrow1$, the case $b=1$ is ruled out by
the same argument.  Therefore $0<a\le\alpha<\beta\le b<1$.
Since the contact points $a$ and $b$ are now interior points, the tangency
observation applies to $L$ at both endpoints.  Thus, the slope of $L$ is
$\varphi'(a)$ and also $\varphi'(b)$.

The inequalities $a\le\alpha$ and $\beta\le b$ are in fact strict.  Suppose
$a=\alpha$.  Then, the slope of $L$ is $\varphi'(\alpha)$.  Since $\varphi'$ is
strictly decreasing immediately to the right of $\alpha$, for small
$x>\alpha$ we have
\[
  \varphi(x)-L(x)
  =
  \int_\alpha^x(\varphi'(t)-\varphi'(\alpha))\,dt
  <0,
\]
again contradicting $L\le\varphi$.  The same argument at $\beta$ shows
$b>\beta$: if $b=\beta$, the tangent line at $\beta$ lies above $\varphi$ just
to the left of $\beta$.
Thus,
\[
  0<a<\alpha<\beta<b<1.
\]

The affine function $L$ is therefore the line through
$(a,\varphi(a))$ and $(b,\varphi(b))$, and the tangency observation gives
\[
  \varphi'(a)=\varphi'(b),
  \qquad
  \varphi(b)-\varphi(a)=\varphi'(a)(b-a),
\]
which are \eqref{eq:contact-equal-slopes}--\eqref{eq:contact-chord-identity}.

There are no other connected components of
\[
  \{x:\widehat\varphi(x)<\varphi(x)\}.
\]
To see this, any such component would be contained entirely in $(0,a)$ or in
$(b,1)$, hence in one of the two strictly convex outer regions.  Repeating
the two-point representation argument above on such a component would show
that $\widehat\varphi$ is affine there and joins its two endpoint contact
points.  But the chord between two points of a strictly convex graph lies
strictly above the graph between them, contradicting
$\widehat\varphi\le\varphi$.  Hence
$\widehat\varphi$ agrees with $\varphi$ outside $[a,b]$ and is the line segment
$L$ on $[a,b]$.  This proves the existence of a pair satisfying
\eqref{eq:contact-equal-slopes}--\eqref{eq:contact-chord-identity} and the asserted shape of the
convex minorant.

We now prove uniqueness of the contact pair.  More precisely, we must
show that there is at most one pair
\[
  a\in(0,\alpha),\qquad b\in(\beta,1)
\]
satisfying the two equations
\eqref{eq:contact-equal-slopes}--\eqref{eq:contact-chord-identity}.  Since the convex-minorant
construction above has already produced one such pair, this will prove
uniqueness of the contact pair in the lemma.

We first observe that on the right convex region $(\beta,1)$, the derivative
$\varphi'$ is strictly increasing.  Therefore, for a fixed
$a\in(0,\alpha)$, there is at most one $b\in(\beta,1)$ satisfying
\[
  \varphi'(b)=\varphi'(a).
\]
Let $G$ be the restriction of $\varphi'$ to the interval $(\beta,1)$.  Then
$G$ is a $C^1$ strictly increasing function with
\[
  G'(b)=\varphi''(b)>0\qquad(\beta<b<1).
\]
Hence, by the inverse function theorem, $G$ has a $C^1$ inverse on its range.
On the set of $a\in(0,\alpha)$ for which
$\varphi'(a)$ belongs to this range, define
\[
  b(a):=G^{-1}(\varphi'(a)).
\]
This is exactly the unique right contact with the same slope as the left
contact $a$ and is $C^1$ as it is a composition of $C^1$ functions.
The set on which $b(a)$ is defined is an interval, because
$\varphi'$ is strictly increasing on the left convex region as well.

By this uniqueness, every possible common-tangent pair must lie on the curve
$a\mapsto(a,b(a))$.
Along this curve, the remaining condition is that the
tangent line at $a$ actually pass through the point $(b(a),\varphi(b(a)))$.
Equivalently, define
\[
  D(a):=\varphi(b(a))-\varphi(a)-\varphi'(a)(b(a)-a).
\]
A zero of $D$ is exactly a pair satisfying the chord equation
\eqref{eq:contact-chord-identity}.  We now show that $D$ is strictly decreasing wherever
it is defined.  Differentiating both sides of the identity
$\varphi'(b(a))=\varphi'(a)$ with respect to $a$ gives
\[
  \varphi''(b(a))\,b'(a)=\varphi''(a).
\]
Differentiating $D$ therefore gives
\begin{align*}
  D'(a)
  &=
  \varphi'(b(a))b'(a)-\varphi'(a)
  -\varphi''(a)(b(a)-a)
  -\varphi'(a)(b'(a)-1)\\
  &=
  \bigl(\varphi'(b(a))-\varphi'(a)\bigr)b'(a)
  -\varphi''(a)(b(a)-a)\\
  &=
  -\varphi''(a)(b(a)-a)<0.
\end{align*}
In the last line we used the equal-slope relation
$\varphi'(b(a))=\varphi'(a)$, the inequality $b(a)>a$, and the strict
convexity on the left region, which gives $\varphi''(a)>0$.  Therefore $D$ can
vanish at most once.  The convex-minorant construction above gives one zero,
so the pair satisfying \eqref{eq:contact-equal-slopes}--\eqref{eq:contact-chord-identity} is
unique.

Finally we show that the contact points depend analytically on $p$, with
$dx_a/dp>0$ and $dx_b/dp<0$.
Define
\[
  F_1(a,b,p)=\varphi_{p,d}'(b)-\varphi_{p,d}'(a),
\]
\[
  F_2(a,b,p)=
  \varphi_{p,d}(b)-\varphi_{p,d}(a)-\varphi_{p,d}'(a)(b-a).
\]
Set $\mathbf F=(F_1,F_2)$.  This map is jointly analytic in $(a,b,p)$ on
$(0,1)^3$.  At a triple $(a,b,p)$ satisfying $F_1=F_2=0$, the derivatives
with respect to $(a,b)$ are
\[
  D_{(a,b)}\mathbf F=
  \begin{pmatrix}
    -\varphi''(a)&\varphi''(b)\\
    -\varphi''(a)(b-a)&0
  \end{pmatrix}.
\]
To obtain this matrix, note that the $b$-derivative of $F_2$ is
$\varphi'(b)-\varphi'(a)$, which
vanishes because $F_1=0$.  The determinant is positive:
\[
  \varphi''(a)\varphi''(b)(b-a)>0.
\]
By the analytic implicit function theorem, the contact points are locally
analytic functions of $p$.  Since the pair satisfying
\eqref{eq:contact-equal-slopes}--\eqref{eq:contact-chord-identity} is unique for each
$p\in(0,p_*)$, these local analytic solutions agree on overlaps.  Hence, they
define analytic functions $x_a(p)$ and $x_b(p)$ on $(0,p_*)$ such that
$F_1(x_a(p),x_b(p),p)=0$ and
$F_2(x_a(p),x_b(p),p)=0$.
By the chain rule, we have
\begin{align}
  D_{(a,b)}\mathbf F
  \binom{da/dp}{db/dp}
  +
  \binom{\partial_pF_1}{\partial_pF_2}
  =0. \label{eq:contact-derivative-system}
\end{align}

If $z=x^{1/d}$, then
\[
  \partial_p\varphi_{p,d}(x)=\frac{p-z}{p(1-p)},
  \qquad
  \partial_p\varphi_{p,d}'(x)=
  -\frac1{p(1-p)d\,z^{d-1}}.
\]
Writing $u_a=a^{1/d}$ and $u_b=b^{1/d}$, this gives
\[
  \partial_pF_1(a,b,p)
  =
  \frac1{p(1-p)d}\left(u_a^{1-d}-u_b^{1-d}\right)>0,
\]
and
\[
  \partial_pF_2(a,b,p)
  =
  \frac1{p(1-p)}
  \left(\frac{u_b^d-u_a^d}{d u_a^{d-1}}-(u_b-u_a)\right)>0.
\]
The last inequality comes from the strict convexity of $z\mapsto z^d$.
As the second row of \eqref{eq:contact-derivative-system} gives
\[
  -\varphi''(a)(b-a)\frac{da}{dp}+\partial_pF_2=0,
\]
by combining $\varphi''(a)>0$ and $b>a$, we obtain $da/dp>0$.

The same linear system also gives the sign of $db/dp$.  The first row gives
\[
  \varphi''(b)\frac{db}{dp}
  =
  \varphi''(a)\frac{da}{dp}-\partial_pF_1.
\]
Substituting the formula for $da/dp$ from the second row, this becomes
\[
  \varphi''(b)\frac{db}{dp}
  =
  \frac{\partial_pF_2}{b-a}-\partial_pF_1.
\]
Set $\Delta=b-a=u_b^d-u_a^d$.  Then
\[
  \frac{\partial_pF_2}{b-a}
  =
  \frac{\partial_pF_2}{\Delta}
  =
  \frac1{p(1-p)}
  \left(
    \frac1{d u_a^{d-1}}
    -\frac{u_b-u_a}{u_b^d-u_a^d}
  \right),
\]
whereas
\[
  \partial_pF_1
  =
  \frac1{p(1-p)}
  \left(
    \frac1{d u_a^{d-1}}
    -\frac1{d u_b^{d-1}}
  \right).
\]
Subtracting the first equation above from the second, the
$1/(d u_a^{d-1})$ terms cancel and we get
\[
  \partial_pF_1-\frac{\partial_pF_2}{b-a}
  =
  \frac1{p(1-p)}
  \left(
    \frac{u_b-u_a}{u_b^d-u_a^d}
    -\frac1{d u_b^{d-1}}
  \right).
\]
The secant slope of the strictly convex function $z\mapsto z^d$ from $u_a$ to
$u_b$ is strictly smaller than its derivative at the right endpoint $u_b$:
\[
  \frac{u_b^d-u_a^d}{u_b-u_a}<d u_b^{d-1}.
\]
Hence, $\partial_pF_2/(b-a)-\partial_pF_1<0$, and since $\varphi''(b)>0$, we get
$db/dp<0$.
\end{proof}

\begin{proof}[Proof of \Cref{lem:contact-point-limits}]
We first verify that the zeros $u_-(p)$ and $u_+(p)$ of $h_{p,d}$ collapse to
$r_*$ as $p\uparrow p_*$.  At $p=p_*$, Lemma~\ref{lem:convexity-defect} gives
$h_{p_*,d}(r_*)=0$, while $h_{p_*,d}$ is strictly decreasing on $(0,r_*)$ and
strictly increasing on $(r_*,1)$.  Hence
\[
  h_{p_*,d}(z)>0\qquad (z\ne r_*).
\]
Fix $\varepsilon>0$ small enough that $r_*-\varepsilon$ and
$r_*+\varepsilon$ lie in $(0,1)$.  By
continuity in $p$,
\[
  h_{p,d}(r_*-\varepsilon)>0,
  \qquad
  h_{p,d}(r_*+\varepsilon)>0
\]
for all $p<p_*$ sufficiently close to $p_*$.  Since $h_{p,d}(r_*)<0$ for such
$p$, the left zero lies in $(r_*-\varepsilon,r_*)$ and the right zero lies in
$(r_*,r_*+\varepsilon)$.  As $\varepsilon>0$ is arbitrary, we conclude that
$u_-(p)\to r_*$ and $u_+(p)\to r_*$ as $p\uparrow p_*$.

For the contact limit as $p\uparrow p_*$, by \Cref{lem:contact-points},
monotonicity gives $x_a(p)\to A$ and $x_b(p)\to B$ for some $A,B\in[0,1]$.
Since $x_a(p)<u_-(p)^d$ and $x_b(p)>u_+(p)^d$, while
$u_\pm(p)\to r_*$, we have $A\le r_*^d\le B$.
After fixing any $p_0\in(0,p_*)$, monotonicity also gives
\[
  0<x_a(p_0)\le x_a(p),
  \qquad
  x_b(p)\le x_b(p_0)<1
  \qquad (p_0<p<p_*),
\]
so $A>0$ and $B<1$.  The map $(p,x)\mapsto\varphi_{p,d}'(x)$ is continuous
at $(p_*,A)$ and $(p_*,B)$.  We may therefore pass to the limit in the
equal-slope equation~\eqref{eq:contact-equal-slopes} to obtain
\[
  \varphi_{p_*,d}'(A)=\varphi_{p_*,d}'(B).
\]
At $p=p_*$, $\varphi_{p_*,d}''\ge0$ and vanishes only at $r_*^d$, so
$\varphi_{p_*,d}'$ is strictly increasing.  Hence
$A=B=r_*^d$.
Taking positive $d$th roots gives
$u_a(p)\to r_*$ and $u_b(p)\to r_*$ as $p\uparrow p_*$.

It remains to consider $p\downarrow0$.  For every fixed $z\in(0,1)$,
\[
  h_{p,d}(z)
  =
  \frac1{1-z}-(d-1)\log\frac{z(1-p)}{(1-z)p}
  \longrightarrow -\infty.
\]
If $0<\eps<r_*$, then $h_{p,d}(\eps)<0$ for all sufficiently small $p$.  Since
$h_{p,d}$ is strictly decreasing on $(0,r_*)$, the left zero satisfies
$u_-(p)<\eps$.  Hence, $u_-(p)\to0$.  Similarly, if
$0<\eps<1-r_*$, then $h_{p,d}(1-\eps)<0$ for all sufficiently small $p$; since
$h_{p,d}$ is strictly increasing on $(r_*,1)$, the right zero satisfies
$u_+(p)>1-\eps$.  Hence, $u_+(p)\to1$.

The contact ordering
\[
  0<u_a(p)<u_-(p)<u_+(p)<u_b(p)<1
\]
then gives $u_a(p)\to0$ and $u_b(p)\to1$ as $p\downarrow0$.  Finally, the
monotonicity from Lemma~\ref{lem:contact-points} turns these limits into
$u_a(p)\downarrow0$ and $u_b(p)\uparrow1$.
\end{proof}

\subsection{Remaining details for the scalar Lubetzky--Zhao boundary theorem}
\label{app:lz-boundary-details}

\begin{proof}[Proof of \textup{(M1)}]
We first check that $\pc(r)<r$.  We shall use that
\[
  p_*=\frac{d-1}{(d-1)+\exp(d/(d-1))}
  <
  \frac{d-1}{d}
  =r_*.
\]
For $r\in(r_*,1)$, the inequality $\pc(r)<r$ is immediate, since
$\pc(r)<p_*<r_*<r$.  Now fix $r\in(0,r_*)$ and set
$p=\pc(r)$.  Then, $r=u_a(p)$ and $p<p_*<r_*<u_b(p)$.  If $u_a(p)\le p$, then
$u_a(p)^d\le p^d<u_b(p)^d$, so $p^d$ lies on the affine part of the convex
minorant joining the two contact points.  But $J_p(p)=0$ is the strict global
minimum of $J_p$.  Thus, unless one of the contact points has first coordinate
$p^d$, the affine part has positive value at $p^d$, contradicting the fact that
the convex minorant lies below $\varphi_{p,d}(p^d)=J_p(p)=0$.  If one contact
point had first coordinate $p^d$, then the common tangent slope would be
$\varphi_{p,d}'(p^d)=0$; the equal-slope condition would force the other contact
point also to have first coordinate $p^d$, since $\varphi_{p,d}'(x)=0$ only at
$x=p^d$.  This contradicts distinctness of the contacts.  Hence
$p<u_a(p)=r$.

At $p=\pc(r)$, the two contact points have first coordinates $r^d$ and
$\sm(r)^d$.  They are distinct by Lemma~\ref{lem:contact-points}, so
$\sm(r)\ne r$.  This proves (M1).
\end{proof}

\begin{proof}[Proof of \textup{(M3)} and \textup{(M4)}]
Fix $r\in(0,1)\setminus\{r_*\}$.  The affine part of the
convex minorant of $\varphi_{\pc(r),d}$ joins
\[
  \bigl(r^d,J_{\pc(r)}(r)\bigr)
  \quad\text{and}\quad
  \bigl(\sm(r)^d,J_{\pc(r)}(\sm(r))\bigr),
\]
and is tangent at both endpoints.  Hence, this affine line is precisely
$\ell_r$, the tangent line at $r^d$.

On the interval between the two contact points, $\ell_r$ is the affine part of
the convex minorant, so it lies below $\varphi_{\pc(r),d}$.  Outside that
interval, it is the tangent line at the adjacent contact point in a strictly
convex outer region, so it again lies below the graph, with equality only at
that contact.  On the open interval between the contact points, the graph is
strictly above the affine part.  Therefore $\ell_r\le\varphi_{\pc(r),d}$ on
$[0,1]$, and the zeros of
\[
  x\mapsto \varphi_{\pc(r),d}(x)-\ell_r(x)
\]
are exactly $r^d$ and $\sm(r)^d$.  This proves (M3) and (M4).
\end{proof}
